IVF partitions the database around coarse centroids and stores a list of
vectors assigned to each centroid. A query chooses $n_{\rm probe}$ lists,
producing a routed universe $U(q)$ of $n_c$ candidates. Visiting fewer lists
reduces access work but can exclude a true neighbour before any fine scoring.
The centroids are a routing device; their distance ordering is a different
selection problem from choosing database vectors.

Product quantisation divides a vector into $M$ coordinate subspaces and
represents each with a label into a small learned codebook \cite{pq}.
If $\widehat x_m=\mathcal C_m[z_m]$, asymmetric distance computation (ADC)
keeps the query unquantised and scores
\begin{equation}
 \|q-\widehat x\|^2=\sum_{m=0}^{M-1}\|q_m-\mathcal C_m[z_m]\|^2.
 \label{eq:2-adc}
\end{equation}
A query-specific table precomputes each subspace distance. Scanning then
uses code loads, table lookups and accumulation. The product is between
subspaces; the centroids \emph{within} a learned subcodebook need not have a
further coordinate product. This distinction determines whether Chapter 5's
exact internal table factorisation is valid.

``Residual'' has two meanings in this setting. IVF-PQ commonly encodes
$x-\mu_{\ell(x)}$, making the coarse centroid part of reconstruction.
JHQ instead routes vectors through IVF while defining its residual relative
to the primary quantised reconstruction in rotated coordinates. These choices
have different codebook dependencies; adding an IVF centroid to JHQ's
residual reconstruction would count a term that its encoder did not subtract.
The original PQ/ADC framework and GPU implementations also permit candidate
expansion followed by a more accurate refinement stage
\cite{johnson_gpu,faiss_refine}.

In the analysed regime with at least $k$ valid routed candidates, JHQ's
primary selection retains valid survivors $C_\alpha(q)$, capped by
$c(\alpha)=\max\{k,\lceil\alpha k\rceil\}$, and refined selection returns
$\mathcal R_\alpha(q)$. The set inclusion is
\begin{equation}
 \mathcal R_\alpha(q)\subseteq C_\alpha(q)\subseteq U(q)\subseteq\mathcal X.
 \label{eq:2-cascade}
\end{equation}
Intersecting with GT gives stage-wise coverage upper bounds on final recall.
These inclusions do not quantify how much loss each stage causes; that
requires fixed-universe counterfactual scoring. Refinement can improve
ordering among retained vectors, but cannot restore excluded IDs.

\par\medskip\noindent\begin{minipage}{\linewidth}
\centering\captionsetup{hypcap=false}
\resizebox{\linewidth}{!}{\begin{tikzpicture}[
 font=\small\linespread{1}\selectfont,
 box/.style={draw=ThesisNavy!65,rounded corners=2pt,align=center,
 text width=2.65cm,minimum height=1.8cm,inner sep=4pt},
 arr/.style={-{Latex[length=2mm]},draw=ThesisNavy!80,thick}]
\node[box] (x) at (0,0) {Database $\mathcal X$\\$N$ vectors};
\node[box] (u) at (3.6,0) {IVF lists\\$U(q)$: $n_c$ vectors};
\node[box,fill=ThesisNavy!7] (c) at (7.2,0) {Primary top-$c$\\$C_\alpha(q)$\\up to $c$ vectors};
\node[box,fill=ThesisNavy!7] (r) at (10.8,0) {Refined top-$k$\\$\mathcal R_\alpha(q)$};
\draw[arr] (x) -- node[above,font=\footnotesize] {IVF} (u);
\draw[arr] (u) -- node[above,font=\footnotesize] {$D_{\rm P}$} (c);
\draw[arr] (c) -- node[above,font=\footnotesize] {$D_{\rm H}$} (r);
\node[align=center,text width=6.4cm] at (2.6,-2.0)
 {Query $q'=\Pi q$ stays FP32.\\Primary codes score the broad scan.};
\node[align=center,text width=5.8cm] at (9.7,-2.0)
 {Residual $r=\Pi x-\widehat p$.\\Survivors decode $\widehat p+\widehat r$.};
\end{tikzpicture}
}
\captionof{figure}[Candidate and precision cascade]{Candidate and precision cascade. IVF selects lists;
primary scoring selects vectors; only survivors read residuals. The hierarchy
is inherited. This diagram defines logical sets and score fidelity,
not measured recall losses or GPU kernel boundaries.}
\label{fig:2-cascade}
\end{minipage}\par\medskip

Ignoring routing and selection, table-based broad scoring costs
$O(n_cM)$ lookups and residual scoring costs $O(cd)$ coordinate work.
This asymmetry motivates selective precision, but does not predict latency:
memory placement, arithmetic, parallel work and candidate selection also
matter. Increasing probes changes $U$; increasing $\alpha$ changes survivors
within that universe. Treating these as one generic search-depth parameter
hides their different quality and cost effects.
